\honest{Correspondence Bias}{Eliezer Yudkowsky}{2007}

I open with Gilbert and Malone: the correspondence bias is ``the tendency to draw inferences about a person's unique and enduring dispositions from behaviors that can be entirely explained by the situations in which they occur.''

When someone else kicks a vending machine, we think they are ``an angry person.'' When we kick it, it is because the bus was late and the machine ate our lunch money twice. We explain our own actions by our situations and other people's by their dispositions, because we do not see their history. \nb{A meta-analysis of 173 studies, published in November 2006, found this self-other difference close to zero on average. It held for negative events, like the kick, and reversed for positive ones.}

Consider the prior probabilities. There are more late buses than people born with unnaturally high anger. \nb{``An angry person'' need not mean an extreme case; someone more irritable than most is common, and against that the late bus does not clearly win.} Everyone carries a few mutations, if I recall correctly, but any given stretch of DNA is unlikely to be affected, and any given trait is probably near average.

Even when people are told that a speaker for or against abortion was randomly assigned the side, they still think the speaker leans that way. \nb{In the experiment cited, by Jones and Harris (1967), the essays were for or against Fidel Castro. As Gilbert and Malone report it, observers did infer attitudes from assigned essays, but much weaker ones than from essays freely chosen.}

It is intuitive to explain rain by water spirits, fire by a fire-stuff escaping from what burns, and sleep by a ``dormitive potency.'' The real mechanisms are evaporation and condensation, oxidation, and chemistry acting on the nervous system; they are harder to think of than essences. So we give the kicker an innate kicking tendency. \nb{Gilbert and Malone, whom the post cites, list four causes of the bias, not this one, and say there is still no accepted explanation.} When we kick the machine ourselves, we also overestimate how many others would do the same: the false consensus effect. Students who drink overestimate how many fellow students drink, and nondrinkers underestimate it. The fundamental attribution error, as I define it, is overrating others' dispositions ``while reversing this tendency for ourselves.'' \nb{The second half is a separate idea, the actor-observer hypothesis of the meta-analysis above.}

The lesson: ``everyone sees themselves as behaving normally.'' Ask what situation people think they are in. People do have dispositions, but there are not enough heritable quirks to account directly for all the behavior we see.

If a red button destroyed the world and a green one stopped it, you would press the green one. Yet people ask me why I want to save the world, as if I had a traumatic childhood. My beliefs are unusual and need explaining, I grant; my reaction does not. Perhaps I overestimate how many would press the green button, but ``I'd still bet there'd be at least a substantial minority.''

Even people who do terrible things are not ``exceptional mutants.'' Once you understand that, you will stop being surprised by human events. \nb{The next day's post, ``Are Your Enemies Innately Evil?'', makes an exception for ``true psychopaths.''}
